% !LW recipe=latexmk (latexmkrc)
% !TEX encoding = UTF-8 Unicode
% !Mode:: "TeX:UTF-8"

\documentclass{styles/resume}
\usepackage{styles/zh_CN-Adobefonts_external}
\usepackage{styles/linespacing_fix}
\usepackage{cite}
\usepackage{titlesec}
\usepackage{enumitem}
\usepackage{graphicx} % 添加图片支持
% 调整 \section 间距：{左缩进}{标题前间距}{标题后间距}
\titlespacing*{\section}{0pt}{3ex plus 1ex minus .2ex}{1ex plus .2ex}
% 调整 itemize 间距：itemsep=item 之间，parsep=同条内段落之间，topsep=列表上下
\setlist[itemize,1]{itemsep=0.8ex, parsep=0.2ex, topsep=0.3ex}
\setlist[itemize,2]{itemsep=0.4ex, parsep=0.1ex, topsep=0.2ex}
\renewcommand{\CJKglue}{\hskip 0.06em}

\begin{document}
\pagenumbering{gobble}

\name{侯汶政｜后端/全栈/Agent工程师}

\section{\faInfoCircle\ 基本信息}
\vspace{0.5ex}
\begin{minipage}[t]{0.48\textwidth}
\begin{itemize}[label=,leftmargin=*]
\large
\item \textbf{当前状态：}\textbf{深圳在职}
\item \textbf{到岗时间：}\textbf{2周左右}
\item \textbf{学历：}\textbf{本科全日制（学信网可查）}
\item \textbf{求职偏好：}不外包 / 不驻场
\end{itemize}
\end{minipage}\hfill
\begin{minipage}[t]{0.48\textwidth}
\begin{itemize}[label=,leftmargin=*]
\large
\item \textbf{微信：}1489938120
\item \textbf{邮箱：}\href{mailto:Soofjan1489938120@gmail.com}{Soofjan1489938120@gmail.com}
\item \textbf{个人网站：}\href{https://Soofjan.com}{Soofjan.com}
\item \textbf{博客：}\href{https://juejin.cn/user/4074147977125020}{技术博客}
\end{itemize}
\end{minipage}
% \begin{minipage}{0.25\textwidth}
% \centering
% {\setlength{\fboxsep}{0pt}\fbox{\includegraphics[height=5cm,keepaspectratio]{photo.jpg}}}
% \end{minipage}

\section{\faCogs\ 综合技术能力}
\begin{itemize}
    \item \textbf{编程语言}
    \begin{itemize}
        \item 深入理解 Go runtime GMP 模型，掌握调度循环、工作窃取、阻塞处理及 goroutine 抢占。
        \item 掌握 Go 并发编程与源码机制，熟悉 channel、slice、map 等组件的实现原理，理解 select、锁及 context 等核心特性。
        \item 深入理解 Go 内存管理与性能排查方法，掌握逃逸和泄漏分析、内存分配器与 GC 原理。
        \item 拥有 Vue/React/Android 实战经验，具备多端适应能力。
    \end{itemize}

    \item \textbf{数据库与中间件}
    \begin{itemize}
        \item 掌握 MySQL InnoDB 存储引擎的 Buffer Pool，事务，MVCC 和锁，索引，内部日志底层以及慢查询的工作原理。
        \item 掌握 Redis 数据结构与底层实现，熟悉过期删除、持久化等机制，掌握缓存三兄弟问题、双写一致性的解决方案。
        \item 具备基础 RabbitMQ 消息队列的知识，理解消息重复、丢失处理方法；了解 Kafka、RocketMQ 相关选型。
        \item 基本了解 Go 全文检索库 Bleve 的索引结构，理解倒排索引的组织方式和查询的基本流程。
    \end{itemize}
    
    \item \textbf{Agent}
    \begin{itemize}
        \item 具备基础的 Agent 的工具调用、状态管理与上下文组织的知识，能够使用 LangGraph 编排条件分支、检查点恢复及人工审核流程。
        \item 基本理解 RAG 的文档处理、检索与生成流程，了解混合检索、结果融合及 Recall@K、MRR 等检索评估方法。
    \end{itemize}

    \item \textbf{其他}
    \begin{itemize}
        \item 熟悉 HTTP/HTTPS、TCP、DNS；结合业务做过 NAT 穿透信令，以及 gRPC + Protobuf。
        \item 熟悉 Docker 容器化部署技术；实现前后端，转发 Docker Engine API 实现镜像、容器、网络和运行状态管理。
        \item 热爱编程与开源，熟练使用 Codex/Cursor 等 AI 工具优化工作效率。
        \item 坚持技术分享与社区建设，发表教学视频与文章 200+，累计 1k+ 点赞收藏。
    \end{itemize}
\end{itemize}

\section{\faGraduationCap\ 教育背景}
\datedsubsection{\textbf{广州大学} \quad 学士，计算机科学与技术}{2019.09 -- 2023.06}
\begin{itemize}
\item CET-4 英语四级（610分），CET-6 英语六级（602分）
\item 全国大学生算法设计与编程挑战赛（银奖）
\end{itemize}

\newpage
\section{\faSuitcase\ 工作经历}
\datedsubsection{\textbf{软件工程师 @ 慧为智能科技有限公司}}{2023.05 - 至今}

\paragraph*{私有云存储系统（NAS）开发}\mbox{}\\
{\small
\noindent NAS（Network Attached Storage，网络附加存储）是部署于局域网内的私有云存储设备，为用户提供文件存储、备份及远程访问等能力。\\
\noindent\textcolor{gray}{\textbf{技术栈：} Go/Gin, Gorm, MySQL, SQLite, Bleve, gRPC, RKNN, Python/FastAPI, PostgreSQL, LangGraph}
}
\vspace{1ex}

\begin{itemize}

    \item \textbf{文档中心}：
    {\small
    \noindent\textcolor{gray}{NAS 本地文档管理与检索服务，支持多格式内容解析、全文索引与检索}
    }
    \vspace{0.2ex}
    \begin{itemize}
        \item 使用多份真实文档对比不同解析方案，使用 pdftotext、antiword 等工具代替LibreOffice，平均内存、cpu降低 56 \%、47 \%。
        \item 针对内存受限环境，基于 Bleve 构建轻量级搜索能力；对超长文档采用分块索引和结果聚合，峰值 heapInuse 下降约 76\%。
        \item 使用 CJDict 词典分词，并以压缩 Trie longest-match 替代 substring + map 查词，将 Trace 总耗时降低约 48\%、进程 user CPU 降低约 63\%。
    \end{itemize}
    \vspace{1ex}

    \item \textbf{相册}：
    {\small
    \noindent\textcolor{gray}{NAS 本地照片管理服务，支持照片检测、去重与时间轴查询}
    }
    \vspace{0.2ex}
    \begin{itemize}
        \item 基于拉普拉斯方差实现照片模糊检测，采用 Worker 并行计算、串行批量写入 SQLite；在 1000 张照片实验中，耗时由 13.51 s 降至 4.41 s。
        \item 实现重复照片识别，结合宽高比、dHash 与 pHash 判定，并按 SHA1 缓存图像哈希；支持重复组管理、保留图切换与批量清理。
        \item 优化相册时间轴慢查询：预计算拍摄时间缺失时的排序时间，回填历史数据并建立复合索引，使分页排序能够使用索引。
    \end{itemize}
    \vspace{1ex}

     \item \textbf{云端}：
    {\small
    \noindent\textcolor{gray}{NAS 远程访问云端信令服务，负责设备发现、穿透凭证发放与连接协商}
    }
    \vspace{0.2ex}
    \begin{itemize}
        \item 主导设计 NAT 穿透交互架构，规范客户端、服务端、云端及穿透端的通信流程，实现复杂网络下稳定高效访问。
        \item 采用应用层 AES+RSA 和传输层 TLS 的双层加密方案，保障凭证与敏感字段在复杂网络下的传输安全。
    \end{itemize}
    \vspace{1ex}

    \item \textbf{MediaAI}：
    {\small
    \noindent\textcolor{gray}{NAS 本地 AI 运行平台，统一管理模型、NPU 资源及 ASR、TTS、人脸分析等 AI 服务}
    }
    \vspace{0.2ex}
    \begin{itemize}
        \item 基于余弦距离实现增量式人脸聚类，基于 LFW/WIDER Face 验证模型效果，实现不同人误合并率 0\%、同人归类准确率 89\%。
        \item 设计 NPU 租约调度机制，支持优先级、deadline 及心跳超时回收；基于 Protobuf/gRPC over Unix Domain Socket 对本机业务提供接口。
    \end{itemize}
    \vspace{1ex}

    \item \textbf{其他}：
    \begin{itemize}
        \item 售后邮件 Agent：一周内完成 MVP，实现邮件分类、知识检索、回复草稿生成与人工审核流程。
        \item 文件索引同步系统：自研 BFS+Queue 架构替代 fsnotify 监听，海量文件场景下，Go 堆峰值约降 92\%；RSS 峰值比 fsnotify 平均低59\%。
    \end{itemize}

\end{itemize}

\end{document}
